% ============================================================================
%  Chapter 3 — Requirements Analysis and System Architecture
% ============================================================================

\chapter{Requirements Analysis and System Architecture}
\label{chap:planning}

\section{Introduction}

This chapter covers the planning work we did before writing any code. We identify who will use the system, list the functional and non-functional requirements, define the product backlog as user stories, model the system with UML diagrams, and describe the four-layer architecture and the overall deployment setup. We also describe the hardware and software environment used during the project.


% ────────────────────────────────────────────────────────────────────────────
\section{Identification of Actors}
\label{sec:actors}

The system interacts with four types of actors, each with a specific role:

\begin{table}[H]
\centering
\begin{tabularx}{\textwidth}{|l|X|}
\hline
\textbf{Actor} & \textbf{Description} \\
\hline
Developer & Writes infrastructure code in an IDE, gets real-time violation feedback through the VS~Code extension, and uses auto-fix to correct issues. \\
\hline
DevOps Engineer & Submits code for analysis through the dashboard or API, reviews scan reports, sets severity thresholds, and manages CI/CD pipelines. \\
\hline
Security Officer & Defines security policies, reviews compliance reports, monitors violation trends through Grafana, and adjusts blocking thresholds. \\
\hline
Jenkins (automated) & Triggers scans automatically when code is pushed via Poll SCM, receives ALLOW/BLOCK decisions, and enforces pipeline pass/fail gates. \\
\hline
\end{tabularx}
\caption{System actors and their roles.}
\label{tab:actors}
\end{table}


% ────────────────────────────────────────────────────────────────────────────
\section{Requirements Capture}
\label{sec:requirements}

\subsection{Functional Requirements}

The functional requirements were gathered through discussions with the project supervisor and organized by how the system integrates with different tools:

\begin{enumerate}[leftmargin=2cm, label=\textbf{FR-\arabic*}, labelsep=0.5cm]
  \item The system shall accept infrastructure code in Terraform, Dockerfile, and Kubernetes YAML formats through a REST API.
  \item The system shall parse each format using a dedicated parser and normalize the result into a unified data model.
  \item The system shall evaluate the normalized code against all applicable security policies.
  \item The system shall return a binary ALLOW or BLOCK decision based on a configurable severity threshold.
  \item Each violation shall include the rule ID, severity level, affected line number, a description, and a fix recommendation.
  \item The system shall provide an auto-fix endpoint that returns corrected code and a git-style diff.
  \item The system shall save scan results in a SQLite database for later analysis.
  \item The system shall generate downloadable HTML and PDF compliance reports.
  \item The system shall expose Prometheus-compatible metrics for use with Grafana.
  \item The system shall provide a web dashboard for interactive code scanning and result viewing.
  \item A VS Code extension shall scan infrastructure files on save and show violations as inline markers.
  \item The system shall integrate with Jenkins pipelines to enforce security gates before deployment.
  \item The system shall support diff-only mode, scanning only lines that changed in a unified diff.
\end{enumerate}

\subsection{Non-Functional Requirements}

\begin{enumerate}[leftmargin=1.5cm, label=\textbf{NFR-\arabic*}]
  \item \textbf{Performance}: Average scan response time must stay under 100\,ms for typical files.
  \item \textbf{Reliability}: The service must maintain full test coverage on critical paths (32/32 tests passing, 0 warnings).
  \item \textbf{Portability}: The entire system must run offline without needing any cloud services.
  \item \textbf{Extensibility}: Adding a new security policy should only require writing a single Python class, without touching existing code.
  \item \textbf{Usability}: Violation messages should be educational --- not just saying what is wrong, but also how to fix it.
  \item \textbf{Maintainability}: The code should follow clean architecture principles with clear separation between layers.
\end{enumerate}


% ────────────────────────────────────────────────────────────────────────────
\section{Product Backlog}
\label{sec:backlog}

The product backlog is organized as user stories in the standard format: \emph{``As a [role], I want [feature] so that [benefit].''}  Table~\ref{tab:product-backlog} shows the complete backlog with priorities.

\begin{table}[H]
\centering
\small
\begin{tabularx}{\textwidth}{|c|X|c|c|}
\hline
\textbf{ID} & \textbf{User Story} & \textbf{Priority} & \textbf{Sprint} \\
\hline
US-01 & As a DevOps engineer, I want to submit Terraform code for security analysis so that I can find misconfigurations before deployment. & High & 1 \\
\hline
US-02 & As a DevOps engineer, I want to submit Dockerfiles for security analysis so that I can spot container security risks. & High & 1 \\
\hline
US-03 & As a DevOps engineer, I want to submit Kubernetes YAML for security analysis so that I can check pod and deployment configurations. & High & 1 \\
\hline
US-04 & As a security officer, I want the system to detect hardcoded secrets so that credentials are never exposed in infrastructure code. & High & 1 \\
\hline
US-05 & As a security officer, I want the system to detect publicly accessible S3 buckets so that data exposure is prevented. & High & 1 \\
\hline
US-06 & As a security officer, I want the system to detect insecure security groups so that unrestricted network access is flagged. & High & 1 \\
\hline
US-07 & As a DevOps engineer, I want an ALLOW/BLOCK decision based on severity so that I can quickly see if the code is ready to deploy. & High & 1 \\
\hline
US-08 & As a security officer, I want comprehensive AWS policies (S3, EC2, RDS, IAM, VPC, CloudTrail) so that cloud resources are thoroughly checked. & High & 2 \\
\hline
US-09 & As a security officer, I want extended Docker and Kubernetes policies so that container configurations are fully covered. & High & 2 \\
\hline
US-10 & As a DevOps engineer, I want a configurable severity threshold so that I can control which violations trigger a BLOCK. & Medium & 2 \\
\hline
US-11 & As a DevOps engineer, I want to scan only changed lines (diff mode) so that existing code does not create noise. & Medium & 2 \\
\hline
US-12 & As a DevOps engineer, I want a web dashboard to scan code interactively so that I can see results without using the command line. & Medium & 3 \\
\hline
US-13 & As a DevOps engineer, I want automatic fixing of common vulnerabilities so that I can fix issues with one click. & High & 3 \\
\hline
US-14 & As a security officer, I want HTML and PDF compliance reports so that I can share audit results with stakeholders. & Medium & 3 \\
\hline
US-15 & As a DevOps engineer, I want scan history stored permanently so that I can track security posture over time. & Low & 3 \\
\hline
US-16 & As a DevOps engineer, I want CI/CD pipeline integration with Jenkins so that insecure code is automatically blocked before deployment. & High & 4 \\
\hline
US-17 & As a developer, I want real-time security scanning in VS Code so that I get feedback before committing code. & High & 4 \\
\hline
US-18 & As an operations engineer, I want Prometheus metrics and Grafana dashboards so that I can monitor security posture in real time. & Medium & 4 \\
\hline
\end{tabularx}
\caption{Product backlog --- user stories.}
\label{tab:product-backlog}
\end{table}


% ────────────────────────────────────────────────────────────────────────────
\section{Global Use Case Diagram}
\label{sec:usecase}

The global use case diagram shows all the interactions between the four actors and the main features of the system. It gives a high-level view of what the microservice does.

\begin{figure}[H]
\centering
\begin{tikzpicture}[
    usecase/.style = {
        draw, thick, ellipse,
        minimum width=2.8cm, minimum height=0.9cm,
        align=center, font=\scriptsize
    },
    arr/.style = {-, thick},
    % Stick figure command via pic
]

% Stick figure helper
\tikzset{
  actor pic/.pic={
    \draw[thick] (0,0.3) circle (0.18cm);
    \draw[thick] (0,0.12) -- (0,-0.3);
    \draw[thick] (-0.25,0) -- (0.25,0);
    \draw[thick] (0,-0.3) -- (-0.18,-0.6);
    \draw[thick] (0,-0.3) -- (0.18,-0.6);
  }
}

% System boundary
\draw[thick, rounded corners=8pt, fill=blue!3]
    (1.5, -7) rectangle (9.5, 1);
\node[font=\small\bfseries, anchor=north west] at (1.7, 0.9) {DevSecOps Policy Service};

% Use cases
\node[usecase, fill=white] (scan)    at (5.5,  0)    {Scan Code};
\node[usecase, fill=white] (fix)     at (5.5, -1.4)  {Auto-Fix\\Violations};
\node[usecase, fill=white] (report)  at (5.5, -2.8)  {Generate\\Report};
\node[usecase, fill=white] (dash)    at (5.5, -4.2)  {View\\Dashboard};
\node[usecase, fill=white] (history) at (5.5, -5.6)  {View Scan\\History};
\node[usecase, fill=white] (metrics) at (5.5, -6.7)  {View Metrics};

% Left actors — Developer
\pic at (-0.8, -0.8) {actor pic};
\node[font=\scriptsize\bfseries, align=center] at (-0.8, -1.6) {Developer};

% Left actors — DevOps Engineer
\pic at (-0.8, -3.5) {actor pic};
\node[font=\scriptsize\bfseries, align=center] at (-0.8, -4.3) {DevOps Eng.};

% Right actors — Security Officer
\pic at (12.2, -1.8) {actor pic};
\node[font=\scriptsize\bfseries, align=center] at (12.2, -2.6) {Security Off.};

% Right actors — Jenkins
\pic at (12.2, -5.4) {actor pic};
\node[font=\scriptsize\bfseries, align=center] at (12.2, -6.2) {Jenkins};

% Developer connections (from feet position ~y=-1.4)
\draw[arr] (-0.8, -1.4) -- (scan);
\draw[arr] (-0.8, -1.4) -- (fix);

% DevOps connections
\draw[arr] (-0.8, -4.0) -- (scan);
\draw[arr] (-0.8, -4.0) -- (fix);
\draw[arr] (-0.8, -4.0) -- (dash);
\draw[arr] (-0.8, -4.0) -- (history);

% Security Officer connections
\draw[arr] (12.2, -2.4) -- (scan);
\draw[arr] (12.2, -2.4) -- (report);
\draw[arr] (12.2, -2.4) -- (metrics);

% Jenkins connections
\draw[arr] (12.2, -6.0) -- (scan);

\end{tikzpicture}
\caption{Global use case diagram showing actors and their interactions with the system.}
\label{fig:global-usecase}
\end{figure}


% ────────────────────────────────────────────────────────────────────────────
\section{Application Architecture}
\label{sec:architecture}

\subsection{Four-Layer Architecture}

The microservice is organized into four clearly separated layers, each with a specific job. This design keeps things clean and makes sure dependencies only flow in one direction~\cite{martin2017cleanarch}.

\begin{table}[H]
\centering
\begin{tabularx}{\textwidth}{|c|l|X|}
\hline
\textbf{Layer} & \textbf{Module} & \textbf{Responsibility} \\
\hline
1 — API & \texttt{api/routes.py} & Handles HTTP requests and responses, validates input using Pydantic schemas, and coordinates the analysis pipeline. \\
\hline
2 — Parser & \texttt{engine/parser.py} & Implements format-specific parsers (Terraform via \texttt{python-hcl2}, Dockerfile via regex, Kubernetes via PyYAML) that turn raw code into structured data. \\
\hline
3 — Normalizer & \texttt{engine/normalizer.py} & Converts the different parser outputs into a single unified \texttt{NormalizedCode} model, which simplifies everything downstream by 67\%. \\
\hline
4 — Policy Engine & \texttt{engine/policy\_engine.py} & Runs all 103 policies against the normalized code, applies the severity threshold, and produces the final ALLOW/BLOCK decision. \\
\hline
\end{tabularx}
\caption{Four-layer architecture of the microservice.}
\label{tab:layers}
\end{table}

The key insight here is the \textbf{Normalizer layer} (Layer~3). It acts as a bridge between the format-specific parsers and the format-agnostic policy engine. Without it, each policy would need to understand every supported format, creating $N \times M$ complexity (where $N$ is the number of policies and $M$ is the number of formats). The normalizer brings this down to $N + M$.

\begin{figure}[H]
\centering
\begin{tikzpicture}[
    layer/.style = {
        draw, thick, rounded corners=4pt,
        minimum width=12cm, minimum height=1.2cm,
        align=center, font=\small\bfseries
    },
    arr/.style = {-{Stealth[length=6pt]}, thick, color=gray!70}
]

\node[layer, fill=blue!15] (l1) at (0,0)
    {Layer 1 --- API \ \ (\texttt{api/routes.py})};
\node[layer, fill=green!15] (l2) at (0,-2)
    {Layer 2 --- Parser \ \ (\texttt{engine/parser.py})};
\node[layer, fill=orange!15] (l3) at (0,-4)
    {Layer 3 --- Normalizer \ \ (\texttt{engine/normalizer.py})};
\node[layer, fill=red!15] (l4) at (0,-6)
    {Layer 4 --- Policy Engine \ \ (\texttt{engine/policy\_engine.py})};

\draw[arr] (l1) -- node[right, font=\scriptsize] {parsed request} (l2);
\draw[arr] (l2) -- node[right, font=\scriptsize] {structured dict} (l3);
\draw[arr] (l3) -- node[right, font=\scriptsize] {NormalizedCode} (l4);

% Return arrow
\draw[arr, dashed, color=blue!60] (l4.west) -- ++(-1.5,0) |- node[left, font=\scriptsize, pos=0.25] {ALLOW/BLOCK} (l1.west);

% Format labels on Parser layer
\node[font=\tiny, anchor=north] at (-3.5,-2.7) {TerraformParser};
\node[font=\tiny, anchor=north] at (0,-2.7) {DockerfileParser};
\node[font=\tiny, anchor=north] at (3.5,-2.7) {KubernetesParser};

\end{tikzpicture}
\caption{Four-layer architecture of the microservice with data flow.}
\label{fig:four-layer-arch}
\end{figure}

\subsection{Global System Architecture}

The complete system is made up of several services working together. Figure~\ref{fig:global-arch} shows the overall deployment setup.

\begin{figure}[H]
\centering
\begin{tikzpicture}[
    component/.style = {
        draw, thick, rounded corners=3pt,
        minimum width=2.8cm, minimum height=1cm,
        align=center, font=\small\bfseries
    },
    arr/.style = {-{Stealth[length=5pt]}, thick},
    darr/.style = {{Stealth[length=5pt]}-{Stealth[length=5pt]}, thick}
]

% Core service
\node[component, fill=blue!20] (api) at (0,0) {Policy Service\\\tiny localhost:8000};

% Clients on the left
\node[component, fill=purple!15] (vscode) at (-5, 2) {VS Code\\Extension};
\node[component, fill=cyan!15] (dashboard) at (-5, 0) {Web\\Dashboard};
\node[component, fill=yellow!15] (jenkins) at (-5, -2) {Jenkins\\\tiny localhost:8080};

% Backend on the right
\node[component, fill=gray!15] (sqlite) at (5, 0) {SQLite\\Database};
\node[component, fill=orange!15] (prom) at (5, -2) {Prometheus\\\tiny localhost:9090};
\node[component, fill=green!15] (grafana) at (5, -4) {Grafana\\\tiny localhost:3001};
\node[component, fill=red!10] (gitea) at (-5, -4) {Gitea\\\tiny localhost:3000};

% Arrows
\draw[darr] (vscode) -- node[above, font=\tiny] {HTTP} (api);
\draw[darr] (dashboard) -- node[above, font=\tiny] {HTTP} (api);
\draw[darr] (jenkins) -- node[above, font=\tiny] {HTTP} (api);
\draw[arr] (api) -- node[above, font=\tiny] {store} (sqlite);
\draw[arr] (prom) -- node[left, font=\tiny] {scrape} (api);
\draw[arr] (grafana) -- node[right, font=\tiny] {query} (prom);
\draw[arr] (jenkins) -- node[left, font=\tiny] {Poll SCM} (gitea);

\end{tikzpicture}
\caption{Global system architecture showing all integrated services.}
\label{fig:global-arch}
\end{figure}

The main components and how they connect:

\begin{itemize}[leftmargin=1.5cm]
  \item \textbf{Policy Service} (\texttt{localhost:8000}): The core FastAPI application that serves both the REST API and the web dashboard.
  \item \textbf{SQLite Database}: An embedded database that stores scan history, violation details, and statistics.
  \item \textbf{VS Code Extension}: A client-side extension that talks to the policy service over HTTP.
  \item \textbf{Jenkins} (\texttt{localhost:8080}): An automation server that runs CI/CD pipelines calling the policy service API.
  \item \textbf{Gitea} (\texttt{localhost:3000}): A self-hosted Git service that provides version control without depending on external services.
  \item \textbf{Prometheus} (\texttt{localhost:9090}): Collects metrics from the \texttt{/metrics/prometheus} endpoint every 15~seconds.
  \item \textbf{Grafana} (\texttt{localhost:3001}): Displays Prometheus metrics through preconfigured dashboards.
\end{itemize}


% ────────────────────────────────────────────────────────────────────────────
\section{Class Diagram}
\label{sec:classdiagram}

The class diagram shows the main classes, their attributes, methods, and relationships across the four layers.

\begin{figure}[H]
\centering
\begin{tikzpicture}[
    class/.style = {
        draw, thick, rectangle,
        minimum width=3.2cm,
        align=center, font=\scriptsize,
        inner sep=3pt
    },
    arr/.style = {-{Stealth[length=5pt]}, thick},
    inh/.style = {-{Triangle[open, length=6pt]}, thick},
    comp/.style = {-{Diamond[open, length=6pt]}, thick}
]

% BasePolicy (abstract)
\node[class, fill=red!10] (base) at (0, 0) {
    \textbf{\guillemotleft abstract\guillemotright}\\\textbf{BasePolicy}\\
    \hrulefill\\
    +rule\_id: str\\
    +severity: PolicySeverity\\
    \hrulefill\\
    +check(code): list
};

% Concrete policies
\node[class, fill=red!5] (p1) at (-3.5, -3) {PublicS3\\BucketPolicy};
\node[class, fill=red!5] (p2) at (0, -3) {Hardcoded\\SecretPolicy};
\node[class, fill=red!5] (p3) at (3.5, -3) {\dots\\ (103 total)};

\draw[inh] (p1) -- (base);
\draw[inh] (p2) -- (base);
\draw[inh] (p3) -- (base);

% PolicyEngine
\node[class, fill=blue!10] (engine) at (7, 0) {
    \textbf{PolicyEngine}\\
    \hrulefill\\
    -policies: list\\
    \hrulefill\\
    +analyze(code, threshold): Result
};

\draw[comp] (engine) -- node[above, font=\tiny] {1..*} (base);

% Parsers
\node[class, fill=green!10] (tf) at (-4, 3.5) {Terraform\\Parser};
\node[class, fill=green!10] (dk) at (0, 3.5) {Dockerfile\\Parser};
\node[class, fill=green!10] (k8) at (4, 3.5) {Kubernetes\\Parser};

% Normalizer
\node[class, fill=orange!10] (norm) at (0, 5.5) {
    \textbf{Normalizer}\\
    \hrulefill\\
    +normalize(parsed): NormalizedCode
};

% NormalizedCode
\node[class, fill=yellow!10] (nc) at (7, 3.5) {
    \textbf{NormalizedCode}\\
    \hrulefill\\
    +code\_type: str\\
    +resources: list[Resource]
};

% AutoFixer
\node[class, fill=purple!10] (fixer) at (7, 5.5) {
    \textbf{AutoFixer}\\
    \hrulefill\\
    +fix(type, code): FixResult
};

% Arrows
\draw[arr] (norm) -- (tf);
\draw[arr] (norm) -- (dk);
\draw[arr] (norm) -- (k8);
\draw[arr] (norm) -- (nc);
\draw[arr] (engine) -- (nc);

\end{tikzpicture}
\caption{Global class diagram of the microservice showing the main classes and their relationships.}
\label{fig:class-diagram}
\end{figure}


% ────────────────────────────────────────────────────────────────────────────
\section{Sprint Planning}
\label{sec:sprint-planning}

Based on the product backlog and the architecture, we distributed the user stories across four sprints. Table~\ref{tab:sprint-planning} shows the detailed allocation.

\begin{table}[H]
\centering
\begin{tabularx}{\textwidth}{|c|X|c|X|}
\hline
\textbf{Sprint} & \textbf{Theme} & \textbf{User Stories} & \textbf{Key Deliverables} \\
\hline
1 & Core Microservice & US-01 to US-07 & FastAPI REST API, 3 parsers, normalizer, 12 base policies, automated tests. \\
\hline
2 & Policy Expansion & US-08 to US-11 & 91 additional policies (total 103), severity threshold, diff-only mode. \\
\hline
3 & Dashboard \& Remediation & US-12 to US-15 & Web dashboard, auto-fix engine, HTML/PDF reports, SQLite history. \\
\hline
4 & Integration \& Monitoring & US-16 to US-18 & Jenkins CI/CD, Gitea, VS Code extension, Prometheus, Grafana. \\
\hline
\end{tabularx}
\caption{Sprint planning --- user story allocation.}
\label{tab:sprint-planning}
\end{table}


% ────────────────────────────────────────────────────────────────────────────
\section{Development Environment}
\label{sec:environment}

\subsection{Hardware Environment}

The project was developed and tested on the following hardware:

\begin{table}[H]
\centering
\begin{tabularx}{\textwidth}{|l|X|}
\hline
\textbf{Component} & \textbf{Specification} \\
\hline
\multicolumn{2}{|l|}{\textbf{Host Machine (Windows)}} \\
\hline
Processor & AMD Ryzen 7 \\
\hline
RAM & 16\,GB DDR4 \\
\hline
Storage & 512\,GB SSD \\
\hline
Operating System & Windows 11 \\
\hline
\multicolumn{2}{|l|}{\textbf{Development VM (Ubuntu 22.04 LTS)}} \\
\hline
vCPUs & 4 cores \\
\hline
RAM & 8\,GB \\
\hline
Storage & 50\,GB virtual disk \\
\hline
OS & Ubuntu 22.04 LTS (64-bit) \\
\hline
\end{tabularx}
\caption{Hardware environment.}
\label{tab:hardware}
\end{table}

\subsection{Software Environment}

\begin{table}[H]
\centering
\begin{tabularx}{\textwidth}{|l|l|X|}
\hline
\textbf{Technology} & \textbf{Version} & \textbf{Purpose} \\
\hline
Python         & 3.10       & Main programming language for the microservice. \\
\hline
FastAPI        & 0.104.1    & Web framework for the REST API~\cite{fastapi2024docs}. \\
\hline
Pydantic       & 2.5.0      & Data validation for request/response schemas~\cite{pydantic2024docs}. \\
\hline
Uvicorn        & 0.24.0     & ASGI server for running the FastAPI application. \\
\hline
python-hcl2    & 4.3.2      & Parser for Terraform HCL files. \\
\hline
PyYAML         & 6.0.1      & Parser for Kubernetes YAML files. \\
\hline
SQLite         & 3.x        & Embedded database for scan history. \\
\hline
Node.js        & 18.x       & Runtime for the VS Code extension. \\
\hline
Jenkins        & 2.x        & CI/CD automation server~\cite{jenkins2024docs}. \\
\hline
Gitea          & Latest     & Self-hosted Git server~\cite{gitea2024docs}. \\
\hline
Prometheus     & Latest     & Metrics collection and time-series database~\cite{prometheus2024docs}. \\
\hline
Grafana        & Latest     & Metrics visualization and dashboards~\cite{grafana2024docs}. \\
\hline
Docker         & Latest     & Containerization of the microservice and monitoring stack. \\
\hline
Docker Compose & Latest     & Multi-container setup for Prometheus and Grafana. \\
\hline
VS Code        & Latest     & IDE used for development and hosting the extension~\cite{vscode2024api}. \\
\hline
pytest         & 7.4.3      & Testing framework for the automated test suite. \\
\hline
Git / GitHub   & Latest     & Version control and remote repository. \\
\hline
\end{tabularx}
\caption{Software environment and technology stack.}
\label{tab:software}
\end{table}


% ────────────────────────────────────────────────────────────────────────────
\section{Conclusion}

This chapter completed the planning phase of the project. We identified the four actors and their roles, captured thirteen functional and six non-functional requirements, organized eighteen user stories into a prioritized product backlog, and described the four-layer architecture in detail. The sprint planning distributed the work across four sprints covering the full project period from February~15 to May~15,~2026, and we listed all the tools and technologies used. The next four chapters present the implementation of each sprint.
